\documentclass[11pt,a4paper]{article}
\usepackage[margin=2.2cm]{geometry}
\usepackage{amsmath,amssymb,booktabs}
\usepackage[hidelinks]{hyperref}

\title{Optimal Stopping with Memory: Deep Recurrent Policies for American Options under
Rough Volatility and Multi-Market Backtesting}
\author{Anamay Mittal, Vraj Shah, Anshul D Jain, Sachit Kapoor, Mridul Kwatra}
\date{MAC F244 Stochastic Calculus --- draft}

% One file per section; owner in brackets. Compile from this folder: pdflatex main.tex
% (the tables are read from ../results/, generated by the scripts in experiments/).
\begin{document}
\maketitle

% ---------------------------------------------------------------- Introduction (P1)
\section{Introduction}

Pricing an American option means deciding when to exercise it, which makes it an
\emph{optimal stopping} problem. The classical numerical methods (binomial trees,
the least-squares Monte Carlo method of Longstaff and Schwartz~\cite{ls2001} and Deep
Optimal Stopping~\cite{becker2019}) all solve it by backward induction. They assume that
the price process is Markovian, so that today's state is enough to decide.

Real markets are not Markovian in the price alone: volatility is persistent and,
empirically, \emph{rough}~\cite{bfg2016}. When the past matters, backward induction needs
an ever larger state, or it carries approximation errors from one date into the
previous one. Damera Venkata and Bhattacharyya~\cite{ospg2023} proposed the
\emph{Optimal Stopping Policy Gradient} (OSPG), which trains a recurrent network to
choose stopping times directly and needs no backward recursion.

We (i) replicate OSPG on the paper's Markovian option benchmarks, (ii) price American
puts under the rough Bergomi model, where the price is non-Markovian, and (iii) backtest
learned exit-timing rules on market data, comparing each with the classical methods.
                 % (P1) introduction and problem statement
% ---------------------------------------------------------------- Theory (P2)
\section{Optimal stopping and American options}

\paragraph{The problem.}
Let $(S_j)_{j=0}^{H}$ be a price process observed at dates $t_0<\dots<t_H$, with
filtration $(\mathcal F_j)$ generated by its history. Stopping at date $j$ pays a reward
$g_j$ that is known at date $j$; for an American (Bermudan) put under the risk-neutral
measure $\mathbb Q$, $g_j=e^{-rt_j}(K-S_j)^+$. A \emph{stopping time} is a random date
$\tau$ with $\{\tau=j\}\in\mathcal F_j$: the decision to stop uses only the past. The price is
\begin{equation}
V_0=\sup_{\tau\in\mathcal T_{0,H}}\ \mathbb E^{\mathbb Q}\big[g_\tau\big].
\end{equation}

\paragraph{Snell envelope.}
Define backwards
\begin{equation}
U_H=g_H,\qquad U_j=\max\big(g_j,\ \mathbb E[U_{j+1}\mid\mathcal F_j]\big),\quad j=H-1,\dots,0 .
\label{eq:snell}
\end{equation}
The process $(U_j)$ is the smallest supermartingale that dominates $(g_j)$; the stopping
time $\tau^*=\min\{j:U_j=g_j\}$ is optimal and $V_0=U_0$~\cite{shreve2}. The term
$C_j=\mathbb E[U_{j+1}\mid\mathcal F_j]$ is the \emph{continuation value}: one exercises
as soon as the immediate reward is at least the value of waiting.

\paragraph{The Markov case.}
If $S$ is Markov and $g_j=G_j(S_j)$, the continuation value is a function of today's state
only, $C_j=c_j(S_j)$. The binomial tree computes $c_j$ exactly on a lattice; Longstaff and
Schwartz estimate it by regressing realised discounted cash flows on functions of
$S_j$~\cite{ls2001}; Deep Optimal Stopping learns the exercise decision at each date with a
neural network of $S_j$~\cite{becker2019}. When the dynamics have memory (as under rough
volatility), $C_j$ depends on the whole path $(S_0,\dots,S_j)$ and these methods lose
information.

\paragraph{Continuous time: a free-boundary problem.}
When exercise is allowed at all times, a put is exercised when $S_t$ falls below a boundary
$L(t)$. Above it the value solves the Black--Scholes equation
$\partial_tV+rS\partial_SV+\tfrac12\sigma^2S^2\partial_{SS}V-rV=0$; below it $V=K-S$; and at
$S=L(t)$ the value meets the payoff smoothly, $\partial_SV=-1$ (smooth pasting). The
boundary is unknown in advance. For a perpetual put it is constant~\cite{shreve2},
\begin{equation}
L^*=\frac{2rK}{2r+\sigma^2},\qquad V(S)=(K-L^*)\Big(\frac{S}{L^*}\Big)^{-2r/\sigma^2}\ \ (S>L^*);
\end{equation}
with $K=40$, $r=0.06$, $\sigma=0.2$ this gives $L^*=30$ and $V(36)=5.79$.

\paragraph{Solving the finite-maturity problem numerically.}
For a finite maturity there is no closed form, and the free-boundary problem is written as
a linear complementarity problem: with $\mathcal L V=(r-\delta)S\partial_SV+\tfrac12\sigma^2S^2\partial_{SS}V-rV$,
\begin{equation}
\min\big(-\partial_tV-\mathcal LV,\ V-g\big)=0,\qquad V(T,S)=g(S).
\end{equation}
Where waiting is optimal the first term vanishes, which is the Feynman--Kac equation for
the discounted expected payoff; where exercise is optimal $V=g$. We solve it in
$x=\log S$, where the coefficients are constant, with Crank--Nicolson time steps (the first
four fully implicit, to damp the oscillations caused by the kink of the payoff) and impose
$V\ge g$ by a penalty iteration~\cite{forsyth2002}: each tridiagonal system is solved with
a large penalty on the nodes where $V<g$, until those nodes stop changing. A Bermudan
option is priced by solving the European equation between exercise dates and taking
$\max(V,g)$ on each of them. Table~\ref{tab:benchmarks} compares the result with a
5{,}000-step-per-year binomial tree: the two independent methods agree to within $0.0004$
on all twenty cases, so either can serve as the ground truth for the simulation methods.
The early-exercise premium is the gap between the American and European columns, and the
Bermudan price with 50 exercise dates per year is only about $0.007$ below the American
one. For $S_0=36$, $\sigma=0.2$, $T=1$ the computed boundary rises from $L(0)=32.9$ to
$L(T^-)=K=40$ and stays above the perpetual boundary $L^*=30$ throughout: with less time
left, waiting is worth less, so exercise starts at higher prices.

\begin{table}[ht]
\centering\small
\caption{American put, $K=40$, $r=0.06$: European (Black--Scholes), Bermudan (50 exercise
dates per year) and American prices from the binomial tree and the finite-difference
solver of the free-boundary problem.}
\label{tab:benchmarks}
\begin{tabular}{rrr r rr rr}
\toprule
$S_0$ & $\sigma$ & $T$ & European & \multicolumn{2}{c}{Bermudan} & \multicolumn{2}{c}{American}\\
 & & & & tree & PDE & tree & PDE\\
\midrule
% S0 & sigma & T & European & Bermudan tree & Bermudan PDE & American tree & American PDE   (generated by experiments/benchmarks.py)
36 & 0.2 & 1 & 3.844 & 4.478 & 4.478 & 4.487 & 4.487 \\
36 & 0.2 & 2 & 3.763 & 4.840 & 4.840 & 4.848 & 4.848 \\
36 & 0.4 & 1 & 6.711 & 7.101 & 7.101 & 7.109 & 7.109 \\
36 & 0.4 & 2 & 7.700 & 8.507 & 8.507 & 8.514 & 8.514 \\
38 & 0.2 & 1 & 2.852 & 3.250 & 3.250 & 3.257 & 3.257 \\
38 & 0.2 & 2 & 2.991 & 3.745 & 3.745 & 3.751 & 3.751 \\
38 & 0.4 & 1 & 5.834 & 6.148 & 6.148 & 6.155 & 6.154 \\
38 & 0.4 & 2 & 6.979 & 7.668 & 7.668 & 7.675 & 7.675 \\
40 & 0.2 & 1 & 2.066 & 2.314 & 2.314 & 2.320 & 2.319 \\
40 & 0.2 & 2 & 2.356 & 2.885 & 2.884 & 2.890 & 2.890 \\
40 & 0.4 & 1 & 5.060 & 5.312 & 5.312 & 5.318 & 5.318 \\
40 & 0.4 & 2 & 6.326 & 6.917 & 6.917 & 6.923 & 6.923 \\
42 & 0.2 & 1 & 1.465 & 1.617 & 1.617 & 1.621 & 1.621 \\
42 & 0.2 & 2 & 1.841 & 2.212 & 2.212 & 2.217 & 2.217 \\
42 & 0.4 & 1 & 4.379 & 4.583 & 4.582 & 4.588 & 4.588 \\
42 & 0.4 & 2 & 5.736 & 6.244 & 6.244 & 6.250 & 6.250 \\
44 & 0.2 & 1 & 1.017 & 1.110 & 1.110 & 1.113 & 1.113 \\
44 & 0.2 & 2 & 1.429 & 1.690 & 1.690 & 1.693 & 1.693 \\
44 & 0.4 & 1 & 3.783 & 3.948 & 3.948 & 3.953 & 3.953 \\
44 & 0.4 & 2 & 5.202 & 5.641 & 5.641 & 5.647 & 5.647 \\

\bottomrule
\end{tabular}
\end{table}

\paragraph{Optional stopping.}
If $(M_j)$ is a martingale and $\tau$ a bounded stopping time, then
$\mathbb E[M_\tau]=M_0$: no stopping rule can improve on a fair game. Two consequences
guide our experiments: stopping a standard Brownian motion is worth exactly~$0$, so memory
can only help when increments are dependent; and if market prices are close to
martingales, no exit rule can beat holding to the horizon in mean, so any benefit of
memory should appear mainly in risk-adjusted terms.
                % (P2) Snell envelope, free boundary, optional stopping
% ---------------------------------------------------------------- Method (P1)
\section{Method: optimal stopping policy gradients}

\paragraph{Stochastic stopping policy.}
A network with parameters $\theta$ reads the history and outputs, at each date,
$\varphi^\theta_j=\mathbb P(\text{stop at } j\mid\text{not stopped before},\ \mathcal F_j)$,
with $\varphi_H=1$. Because stopping does not change the path, the distribution of the
stopping time on a given path is exact~\cite{ospg2023}:
\begin{equation}
\psi^\theta_j=\mathbb P(\tau=j)=\varphi^\theta_j\prod_{n<j}\big(1-\varphi^\theta_n\big),\qquad
\psi^\theta_H=\prod_{n<H}\big(1-\varphi^\theta_n\big).
\end{equation}
The expected reward $J(\theta)=\mathbb E\big[\sum_{j}\psi^\theta_j g_j\big]$ can therefore be
evaluated on pre-simulated paths, without rollouts and without backward recursion.

\paragraph{Gradient and loss.}
By the log-derivative trick,
\begin{equation}
\nabla_\theta J=\mathbb E\Big[\sum_{j=0}^{H}(g_j-b)\,\psi^\theta_j\,\nabla_\theta\log\psi^\theta_j\Big],
\end{equation}
for any constant $b$, since $\sum_j\psi_j=1$ makes the $b$ term vanish. We minimise the
temporal loss $\mathcal L=-\frac1N\sum_{i,j}\mathrm{sg}\big[\psi_{ij}(g_{ij}-b)\big]\log\psi_{ij}$,
where $\mathrm{sg}$ stops gradients, whose gradient is exactly $-\nabla_\theta J$ (checked
numerically in our tests). The update raises $\varphi_j$ when the ratio of the immediate
value $\psi_jg_j$ to the continuation value $\sum_{n>j}\psi_ng_n$ exceeds the current odds
$\varphi_j/(1-\varphi_j)$, an implicit Snell-envelope comparison at every date.

\paragraph{Architecture and training.}
The recurrent policy (OSPG-GRU) is a single-layer GRU with 20 hidden units followed by a
linear layer and a sigmoid; the memoryless policy (OSPG-MLP) applies a two-layer
network with 20 units to each date's inputs, with batch normalisation of its input and of
each hidden layer's output, as in the paper; at decision time the normalisation uses the
statistics from training, so a path's decision never depends on the other paths. Inputs at date $j$ are the standardised
features, the current reward and $j/H$, all known at date~$j$. We train with Adam
(learning rate $10^{-3}$), keep the parameters with the best validation value, and stop
at the first date with $\varphi_j\ge\frac12$ at test time. Three implementation choices
were needed. The output bias starts at $-2$, because a policy started near
$\varphi=\frac12$ converged to ``stop immediately'' on the in-the-money put. A small
penalty $10^{-3}\,\overline{z^2}$ on the pre-sigmoid outputs keeps $\varphi$ away from $0$:
the policy gradient is proportional to $\varphi$, and without the penalty the GRU policy
stalled at ``never stop'' on the 50-asset max-call. Finally, date~$0$, where all simulated
paths coincide, is decided by comparing $g_0$ with the value achieved by the trained
policy, as for the baselines.

\paragraph{Baselines.}
Longstaff--Schwartz regression on $[1,x,x^2]$ of the standardised features and current
reward (in-the-money paths only for options); Deep Optimal Stopping with one network per
date, warm-started from the next date; and fitted Q-iteration with one network (MLP or
GRU) for all dates~\cite{ospg2023}. The baseline networks are run without batch
normalisation: on the American put it lowered FQI-MLP from $4.480$ and $4.462$ to $4.426$
and $4.442$ (two seeds), so we kept the stronger configuration for the baselines, which
can only make the comparison harder for OSPG. Every method sees the same features, and
all values are measured out of sample, so they are lower bounds on the true price; the
binomial tree and the finite-difference solver of Section~2 give the exact prices they
are measured against.
                % (P1) OSPG derivation and implementation, baselines

\section{Experiments and results}
% ---------------------------------------------------------------- Replication results (P1)
\subsection{Replication of the option benchmarks}
\label{sec:replication}

\paragraph{Set-up.}
We ran three Markovian problems for which the exact price is known.
\emph{E1}: the Bermudan put of Section~2 ($S_0=36$, $K=40$, $r=0.06$, $\sigma=0.2$, $T=1$,
50 exercise dates; exact price $4.478$ from the tree and the PDE solver).
\emph{E2}: the paper's Bermudan max-call on $d=20$ and $50$ independent assets
($s_0=100$, 9 exercise dates, Table~1 of~\cite{ospg2023}); the reference is the midpoint of
the lower and upper bounds of Becker et al.~\cite{becker2019}, $[51.571, 51.765]$ and
$[69.582, 69.889]$.
\emph{E3}: the paper's American call on the geometric average of $d=7$ and $20$ correlated
assets (100 exercise dates, Appendix~D.3); the geometric average of correlated geometric
Brownian motions is itself one, so the exact price follows from a one-dimensional tree.
Each problem was simulated with 100{,}000 paths (40{,}000 for E2), split at random into
50\% test and 50\% training, of which a fifth is held out for validation, and repeated for
three seeds. Every method sees the same features and is scored on the same test paths.

\begin{table}[ht]
\centering\small
\caption{Out-of-sample values, mean over three seeds (standard deviation across seeds).
The standard error of one test set is about $0.013$ (put), $0.16$ (max-call) and $0.065$
(geometric call). DOS was not run on the 100-date geometric call.}
\label{tab:replication}
% generated by experiments/summarize.py: method & mean (sd) per problem
\begin{tabular}{lrrrrr}
\toprule
Method & Put & Max-call $d{=}20$ & Max-call $d{=}50$ & Geo-call $d{=}7$ & Geo-call $d{=}20$ \\
\midrule
LSM & 4.458 (0.004) & 51.69 (0.20) & 69.83 (0.32) & 10.169 (0.012) & 9.899 (0.011) \\
DOS & 4.465 (0.007) & 51.69 (0.16) & 69.72 (0.29) & -- & -- \\
FQI-MLP & 4.467 (0.011) & 51.68 (0.12) & 69.67 (0.21) & 10.145 (0.068) & 9.874 (0.056) \\
FQI-GRU & -- & 51.73 (0.14) & 69.81 (0.37) & -- & -- \\
OSPG-MLP & 4.473 (0.008) & 51.80 (0.20) & 69.93 (0.28) & 10.255 (0.046) & 9.989 (0.017) \\
OSPG-GRU & 4.465 (0.015) & 51.72 (0.21) & 69.87 (0.25) & 10.237 (0.065) & 9.941 (0.030) \\
\midrule
Reference & 4.478 & 51.67 & 69.74 & 10.260 & 10.033 \\
Paper, DNN-OSPG & -- & 51.02 & 69.06 & 10.252 & 10.010 \\
\bottomrule
\end{tabular}

\end{table}

\paragraph{Every method recovers the price; OSPG-MLP comes closest.}
Table~\ref{tab:replication} shows all methods within $1.6\%$ of the exact price. The
memoryless OSPG policy is the most accurate on the three problems where the differences
are larger than the noise: $-0.12\%$ on the put, and $-0.05\%$ and $-0.44\%$ on the
geometric call, against $-0.9\%$ to $-1.6\%$ for Longstaff--Schwartz and fitted
Q-iteration. This is the paper's ranking (OSPG ahead of FQI and LS), and our OSPG errors are
of the same size as the paper's ($0.07\%$ and $0.23\%$), although on the geometric call we
trained on about thirty times fewer paths. On the max-call every method is within one standard error of the
reference. OSPG-MLP's $51.80$ and $69.93$ lie slightly above Becker et al.'s upper bounds,
but a mean of three test sets has a standard error of about $0.09$, so neither is
significant.

\paragraph{Paired comparisons.}
Because all methods are evaluated on the same test paths, we compare them path by path:
the noise common to both methods cancels and much smaller differences become visible
(\texttt{results/paired.md}). OSPG-GRU is ahead of Longstaff--Schwartz and FQI-MLP on both
geometric calls by $0.04$ to $0.09$ ($z$ between $2.6$ and $4.6$), and of DOS and FQI-MLP on
the 50-asset max-call by $0.15$ and $0.20$ ($z=3.9$ and $5.4$). It is, however, behind
OSPG-MLP on the put ($0.008$, $z=3.0$), on both max-calls ($0.086$ and $0.064$, $z=2.9$ and
$2.2$) and on the 20-asset geometric call ($0.049$, $z=3.7$).

\paragraph{Memory has a small cost when there is nothing to remember.}
In these Markovian problems the history carries no information beyond today's state,
so the recurrent policy has to learn to ignore it, and it does so imperfectly: it loses
$0.1\%$ to $0.5\%$ against the memoryless policy and trains 1.3 to 4 times more slowly.
Part of the gap may come from batch normalisation, which the MLP uses (as in the paper)
and the GRU does not. The loss is small next to the margin by which both OSPG policies beat
the backward-induction methods on the harder problems, so the recurrent policy remains a
reasonable default when it is not known whether the problem has memory.

\paragraph{The paper's own max-call values.}
The paper reports $50.86$--$51.02$ for $d=20$ and $68.87$--$69.06$ for $d=50$, about $0.8$
below Becker et al.'s lower bounds and below every method here, including our
implementations of the paper's baselines. We therefore compare with Becker et al.\ rather than with the
paper's numbers; the paper's conclusions about the max-call (all methods within a few
tenths of each other) are unaffected.

% Note for P3 (rough Bergomi, stand-in simulator, results/paired.md): both OSPG policies
% beat LSM, DOS and FQI-GRU by 0.035-0.08 (z = 2.3-4.5), but OSPG-GRU and OSPG-MLP are equal
% (difference 0.004, z = 0.3). On the stand-in, the gain comes from avoiding backward
% induction, not from reading the history. Re-check with your simulator and an H sweep.
   % (P1) replication of the option benchmarks, E1-E3
% P3: fractional Brownian motion and rough volatility
\subsection{Fractional Brownian motion and rough volatility}
\label{sec:roughvol}

\paragraph{Fractional Brownian motion.}
For a Hurst parameter $h\in(0,1)$, fractional Brownian motion (fBm)
$W^h$ is a centred Gaussian process with covariance
\begin{equation}
  \mathbb E[W_t^h W_s^h]
  =\frac12\left(t^{2h}+s^{2h}-|t-s|^{2h}\right).
  \label{eq:fbm-covariance}
\end{equation}
When $h=1/2$, fBm is standard Brownian motion. For $h<1/2$ its increments
are negatively correlated, while for $h>1/2$ they are positively correlated.
For $h\ne1/2$, fBm is not a semimartingale, so the usual Markovian
state representation is not available.

The simulator generates exact finite-dimensional Gaussian samples on the
chosen time grid. It constructs the covariance matrix $C$ from
\eqref{eq:fbm-covariance}, computes its Cholesky factorisation
$C=LL^\top$, draws independent standard normal vectors $Z$, and sets
$W=ZL^\top$. The deterministic initial value $W_0^h=0$ is then prepended
to every path. As a numerical check, for $h=0.3$, $40{,}000$ paths and
10 time intervals, the empirical terminal variance was $1.0022$ against
the theoretical value $1.0000$; the maximum absolute difference between
the empirical and theoretical covariance matrices was $0.00792$.

\paragraph{fBm stopping experiment.}
Following the stopping problem in the OSPG paper~\cite{ospg2023}, the reward is the fBm
value itself and the objective is
\begin{equation}
  \sup_{0\leq\tau\leq H}\mathbb E[W_\tau^h].
\end{equation}
The observation supplied to a policy contains the current fBm value and
normalised time; the reward at date $j$ is $R_j=W_{t_j}^h$. We tested
$h\in\{0.05,0.10,0.20,0.30,0.40,0.50,0.60,0.70,0.80,0.90,0.95\}$
using 20,000 paths, 100 time intervals and two random 50/50 train--test
splits, with 20\% of the training portion reserved for validation. The
methods were DOS, FQI-MLP, FQI-GRU, OSPG-MLP and OSPG-GRU. This is a
reduced-compute reproduction; the paper used 40,000 paths and ten splits.

\begin{table}[ht]
\centering
\scriptsize
\caption{fBm stopping results. Entries are mean test rewards; parentheses
show the standard deviation across the two split-level estimates.}
\label{tab:fbm-results}
\begin{tabular}{rccccc}
\toprule
$h$ & DOS & FQI-MLP & FQI-GRU & OSPG-MLP & OSPG-GRU \\
\midrule
0.05 & 0.700 (0.005) & 0.910 (0.017) & 1.297 (0.001) & 1.153 (0.004) & 1.286 (0.004) \\
0.10 & 0.508 (0.024) & 0.740 (0.004) & 1.040 (0.004) & 0.926 (0.003) & 1.032 (0.008) \\
0.20 & 0.304 (0.002) & 0.481 (0.009) & 0.643 (0.009) & 0.578 (0.004) & 0.639 (0.002) \\
0.30 & 0.187 (0.007) & 0.299 (0.008) & 0.366 (0.005) & 0.334 (0.001) & 0.362 (0.008) \\
0.40 & 0.092 (0.021) & 0.132 (0.008) & 0.158 (0.012) & 0.141 (0.010) & 0.151 (0.015) \\
0.50 & -0.001 (0.007) & -0.002 (0.009) & -0.006 (0.003) & -0.004 (0.006) & -0.008 (0.011) \\
0.60 & 0.070 (0.003) & 0.081 (0.013) & 0.099 (0.010) & 0.096 (0.003) & 0.094 (0.007) \\
0.70 & 0.165 (0.002) & 0.177 (0.000) & 0.202 (0.001) & 0.191 (0.002) & 0.193 (0.002) \\
0.80 & 0.245 (0.005) & 0.258 (0.004) & 0.260 (0.001) & 0.266 (0.004) & 0.262 (0.002) \\
0.90 & 0.302 (0.007) & 0.306 (0.003) & 0.282 (0.013) & 0.314 (0.008) & 0.306 (0.007) \\
0.95 & 0.353 (0.001) & 0.346 (0.004) & 0.328 (0.028) & 0.357 (0.002) & 0.349 (0.001) \\
\bottomrule
\end{tabular}
\end{table}

At $h=0.5$, all five learned-method estimates are close to zero, as
expected for a bounded stopping rule applied to a Brownian martingale.
Away from $h=0.5$, the attainable mean reward is positive. In this
reduced experiment, the recurrent models are competitive, but the GRU
variants do not dominate the MLP variants at every $h$. The results
therefore support the role of increment dependence in the stopping
problem, but do not establish a uniform GRU advantage. The two-split
sample is small, so differences between methods should be interpreted
as descriptive rather than as conclusive statistical rankings.

\paragraph{Rough Bergomi model and hybrid simulation.}
For the rough-volatility experiments we use the rough Bergomi model
\begin{align}
  dS_t &= rS_t\,dt+\sqrt{v_t}S_t\,dZ_t,\\
  Z_t &= \rho W_t+\sqrt{1-\rho^2}\,W_t^\perp,\\
  v_t &= \xi_0\exp\!\left(\eta\widetilde W_t^h
                    -\frac{\eta^2}{2}t^{2h}\right),\\
  \widetilde W_t^h &= \sqrt{2h}\int_0^t(t-s)^{h-1/2}\,dW_s.
  \label{eq:rough-bergomi}
\end{align}
Here $W$ and $W^\perp$ are independent Brownian motions, and $\rho$
controls the correlation between the stock-price shock and the driver
of volatility. The variance normalisation gives
$\mathbb E[v_t]=\xi_0$.

The rough Bergomi specification follows Bayer, Friz and Gatheral~\cite{bfg2016}. The active simulator uses a hybrid discretisation of the Volterra
integral: the near-diagonal Brownian integral is sampled using its
conditional Gaussian law given the corresponding Brownian increment;
the remaining kernel contributions are approximated using the hybrid
scheme's evaluation weights. The previous joint-Cholesky implementation
is retained as \texttt{rbergomi\_paths\_exact} for comparison. Stock prices are
then advanced in log space with a left-point variance step. For the
American put, the discounted reward is
\begin{equation}
  R_j=e^{-rt_j}(K-S_{t_j})^+,
\end{equation}
and the base features at each exercise date are $S_{t_j}/K$ and
$v_{t_j}/\xi_0$.

\paragraph{Simulator validation.}
The hybrid simulator passed the repository's existing 71-test suite;
the additional P3 variance test also passed separately. A preliminary
European-call check used $h=0.07$, $\eta=1.9$, $\rho=-0.9$,
$\xi_0=0.0552$, $S_0=K=1$, $r=0$, 50 time intervals and 30,000 paths.
The Monte Carlo estimate was $0.07886$ with standard error $0.00061$,
compared with the published reference value $0.0791$. The difference is
about 0.39 standard errors. This is a useful consistency check, not a
high-precision benchmark reproduction, because the path count and time
grid are smaller than those used for the reference calculation.

\paragraph{American put experiment (E4).}
We ran a reduced Hurst sweep at $h\in\{0.07,0.10,0.30,0.50\}$ using
3,000 simulated paths, 25 time intervals and one random split per $h$.
The model parameters were $T=1$, $\eta=1.9$, $\rho=-0.9$,
$\xi_0=0.09$, $r=0.05$, $S_0=K=100$. We compared Longstaff--Schwartz
(LSM), Deep Optimal Stopping (DOS), OSPG-MLP and OSPG-GRU. The table
reports the test-set mean reward and its pathwise standard error.

\begin{table}[ht]
\centering
\small
\caption{Reduced E4 American put results under hybrid rough Bergomi.
Entries are mean discounted payoffs (standard error).}
\label{tab:e4-results}
\begin{tabular}{rcccc}
\toprule
$h$ & LSM & DOS & OSPG-MLP & OSPG-GRU \\
\midrule
0.07 & 7.935 (0.340) & 7.985 (0.350) & 8.136 (0.340) & 7.932 (0.356) \\
0.10 & 7.698 (0.329) & 7.776 (0.348) & 7.934 (0.340) & 7.690 (0.359) \\
0.30 & 7.362 (0.337) & 7.435 (0.350) & 7.507 (0.332) & 7.367 (0.368) \\
0.50 & 7.540 (0.314) & 7.619 (0.315) & 7.605 (0.352) & 7.479 (0.360) \\
\bottomrule
\end{tabular}
\end{table}

OSPG-MLP has the largest point estimate in these four pilot runs.
However, the standard errors are around $0.31$--$0.37$, and the
uncertainty intervals overlap substantially. With only one split and
3,000 paths, these results do not establish that OSPG-MLP is superior
to the other methods, nor do they establish a GRU advantage. The sweep
is a preliminary computational comparison; more paths and independent
splits would be needed for firm statistical conclusions.
             % (P3) fBm and rough volatility, E4
% ---------------------------------------------------------------- Data and backtest methodology (P4)
\subsection{Data and backtest methodology}
\label{sec:data}
% P4: markets and sources, cleaning and survivorship, windows, features, walk-forward splits,
% cost models, rule-based exits.
         % (P4) data and backtest methodology
% ---------------------------------------------------------------- Backtest results (P5)
\subsection{Backtest results}
\label{sec:backtest}
% P5: results across asset classes, both rewards, block-bootstrap intervals, regimes,
% cost sensitivity, GRU advantage vs volatility persistence, leakage checks.
      % (P5) backtest results, E5

% ---------------------------------------------------------------- Conclusion (P2)
\section{Conclusion and limitations}

\paragraph{What we found.}
Optimal stopping can be solved without backward induction. The OSPG policy gradient
optimises the expected reward of a stochastic stopping rule directly, using the exact
distribution of its stopping time on pre-simulated paths; it needs neither rollouts nor
the recursive continuation values on which the Snell-envelope methods rely. On the
Markovian option benchmarks the memoryless OSPG policy comes closest to the exact
prices, within $0.12\%$ on the put and $0.44\%$ on the 20-asset geometric call, ahead of
Longstaff--Schwartz and fitted Q-iteration (Section~\ref{sec:replication}); the binomial
tree and our finite-difference solver of the free-boundary problem, which agree with each
other to $4\times10^{-4}$, provide the exact prices. A recurrent policy pays a small price
for memory it does not need ($0.1$--$0.5\%$ on these problems), but still beats the
backward-induction methods where the problems are hard.
% P3: one or two sentences on the rough Bergomi put (does reading the history help?).
% P5: one or two sentences on the backtest (does any rule beat holding, net of costs?).

\paragraph{Limitations.}
\begin{itemize}
\item All values are out of sample and therefore lower bounds; we report no upper bounds
  (dual methods), so the distance to the true price is known only where an exact price exists.
\item For the 100-date geometric call we trained on 40{,}000 paths, about thirty times fewer
  than the paper's 1.28 million, and our errors there are correspondingly larger.
\item OSPG needed three practical fixes (initial bias, logit penalty, separate date-0
  decision) that the paper does not describe; without them it can converge to ``stop
  immediately'' or ``never stop''.
\item The stopping rule is a hard threshold $\varphi_j\ge\frac12$ on a stochastic policy
  trained for the soft objective; the two coincide only when the policy is nearly
  deterministic.
\item The memory results rest on one rough-volatility model and one set of markets.
  % P3/P5: add model-specific limitations.
\end{itemize}
            % (P2) conclusion and limitations

\begin{thebibliography}{9}
\bibitem{ospg2023} N. Damera Venkata and C. Bhattacharyya. Deep recurrent optimal stopping.
  \emph{NeurIPS}, 2023.
\bibitem{becker2019} S. Becker, P. Cheridito and A. Jentzen. Deep optimal stopping.
  \emph{JMLR} 20(74), 2019.
\bibitem{ls2001} F. A. Longstaff and E. S. Schwartz. Valuing American options by simulation:
  a simple least-squares approach. \emph{Review of Financial Studies} 14(1), 2001.
\bibitem{bfg2016} C. Bayer, P. Friz and J. Gatheral. Pricing under rough volatility.
  \emph{Quantitative Finance} 16(6), 2016.
\bibitem{forsyth2002} P. A. Forsyth and K. R. Vetzal. Quadratic convergence for valuing
  American options using a penalty method. \emph{SIAM Journal on Scientific Computing}
  23(6), 2002.
\bibitem{shreve2} S. E. Shreve. \emph{Stochastic Calculus for Finance II: Continuous-Time
  Models}. Springer, 2004.
\end{thebibliography}
\end{document}
